\phantomsection\label{fig:vol01:spring-autumn:zheng-between-jin-chu}
\begin{center}
\begin{tikzpicture}[x=.88cm,y=.88cm,font=\small]
  \fill[paper] (0,0) rectangle (11.2,5.8);

  \draw[source!70,line width=1pt]
    (.4,4.8) .. controls (2.8,4.45) and (5.1,5.15) .. (10.7,4.6);
  \node[text=source,above] at (5.6,4.85) {黄河与中原北缘方向};
  \draw[source!55,line width=.85pt]
    (7.3,.4) .. controls (6.75,1.9) and (7.7,2.4) .. (9.8,3.35);
  \node[text=source,below] at (8.1,1.35) {汉水—北上通道方向};

  \fill[dynasty!19] (.85,3.0) -- (3.25,2.9) -- (3.5,4.65) -- (1.1,4.8) -- cycle;
  \node[align=center] at (2.2,3.85) {晋\\北方联盟方向};
  \fill[timeline!22] (4.5,2.05) -- (6.65,2.0) -- (6.9,3.7) -- (4.75,3.9) -- cycle;
  \node[align=center] at (5.7,2.95) {郑\\近畿枢纽};
  \fill[source!17] (8.1,.75) -- (10.4,.7) -- (10.6,2.55) -- (8.35,2.7) -- cycle;
  \node[align=center] at (9.35,1.65) {楚\\江汉方向};
  \fill[timeline!13] (7.2,3.0) -- (8.4,2.95) -- (8.55,3.85) -- (7.35,3.95) -- cycle;
  \node[align=center] at (7.9,3.45) {陈、蔡等\\中原关系};

  \draw[-{Stealth[length=2mm]},ultra thick,color=source]
    (8.25,1.95) .. controls (7.25,2.2) and (6.9,2.65) .. (6.8,2.8);
  \draw[-{Stealth[length=2mm]},thick,dashed,color=dynasty]
    (3.35,3.65) .. controls (4.15,3.3) and (4.35,3.0) .. (4.6,2.95);
  \draw[{Stealth[length=2mm]}-{Stealth[length=2mm]},thin,dashed,color=timeline]
    (6.85,3.35) -- (7.2,3.4);

  \node[text=source,above,align=center] at (7.2,2.2)
    {前598年：楚围郑\\郑暂时屈楚};
  \node[text=dynasty,above,align=center] at (3.75,3.95)
    {前562年：晋率诸侯伐郑};
  \node[text=timeline,above,align=center,font=\footnotesize] at (7.1,4.0)
    {争取中原关系};
  \node[text=source,align=center,font=\footnotesize] at (5.65,.8)
    {两条箭线相隔数十年，用以显示同一走廊反复受压，\\并非同时行军或可精确复原的道路};
\end{tikzpicture}

\smallskip
\small\textit{图\quad 郑在晋楚竞争之间的走廊示意：以前598年楚围郑和前562年晋率诸侯伐郑为两个时间层，说明郑及周边中原关系何以反复承受征发与外交压力。参考《春秋·宣公十一年、襄公十一年》《左传·宣公十一年、襄公十一年》《史记·楚世家》及《史记·晋世家》；参战国、城内交涉和具体行程多赖传世叙事。各政治体范围、通道和箭线均为约略示意，不表示精确疆界、行军路线或同时存在的联盟关系。}
\end{center}
